% ---------------------------------------------------------------
\section{Limitations and Future Work}
\label{sec:limitations-and-future-work}
% ---------------------------------------------------------------

The Bailout Protocol delivers a compulsory, deterministic escalation path from any spawned node to a named human Purpose Layer anchor, with safety, liveness, and accountability properties established under explicit structural assumptions. Four honest limitations bound the scope of those guarantees; each defines a directed research agenda for subsequent work.

% ---------------------------------------------------------------
\subsection{Limitations}
\label{sec:limitations}

\subsubsection{Human Response Latency}
\label{sec:lim-latency}

The termination theorem (Theorem~\ref{thrm:termination}) guarantees a \emph{finite} escalation path to a human anchor, not a \emph{short} one. Worst-case human response time is a function of time-of-day, geographic distribution, communication channel availability, and cognitive load --- all variables entirely outside the protocol's control. A precision irrigation controller at 03:00\,AM that stalls in \texttt{ESCALATING} state for 45 minutes formally satisfies the $T_{\max}$ bound but may violate operational SLAs for water-volume budgets in life-critical agricultural deployments.

The Training Wheels Protocol (Section~\ref{sec:purpose-anchor}) intercepts outbound actions before execution and reduces reckless autonomy, but it does not accelerate the bailout path itself. High-frequency operational domains --- algorithmic trading, autonomous vehicle navigation, real-time process control --- face a fundamental mismatch between millisecond time constants and human cognitive response bandwidth. The protocol correctly stalls execution rather than proceeding without authorization; the cost of that stall is domain-dependent and must be made explicit in any safety case for life-critical deployments. The $T_{\max}$ bound provides a worst-case envelope, but envelope adequacy is a deployment certification question, not a protocol design question.

\subsubsection{Adversarial Conditions}
\label{sec:lim-adversarial}

The correctness theorems (Theorems~\ref{thrm:drift-prevention} and~\ref{thrm:chain-integrity}) assume the sealed envelope, immutable parent pointer, and forensic ledger resist adversarial modification within a hardware-fault model. Under a strong adversarial model --- one in which an attacker has achieved code execution inside a node's Fact Layer process --- these guarantees degrade.

An adversary who compromises a sealed envelope may suppress the bailout trigger itself, preventing the escalation signal from reaching the Purpose Layer. This attacks the \emph{detection} mechanism rather than the parent pointer chain. The forensic ledger records spawn events, but if the Fact Layer process is compromised before suppression, ledger entries can be written to obscure the attack's痕迹. The system responds conservatively by triggering bailout on anomalous absence rather than explicit alarm, but designing against an active, informed Fact Layer adversary is architecturally orthogonal to this protocol.

The applicability boundary is deployments where the host platform is trusted but individual agents may behave unexpectedly: the standard assumption for production AI in regulated industries. National infrastructure, military systems, and financial clearinghouses require substrate-integrity mechanisms --- secure boot chains, hardware-attested execution environments, independent watchdog monitors on separate hardware --- that compose with but are architecturally outside this protocol's scope.

\subsubsection{Scalability Under 1:1 Anchoring}
\label{sec:lim-scalability}

The termination guarantee (Theorem~\ref{thrm:termination}) bounds hierarchy \emph{depth} but not \emph{breadth}. Under the current 1:1 anchoring model, a single human anchor $p \in P$ may supervise many independent subtrees, each spawning children to depth $d_{\max}$. As active node count grows, spawn events, ledger entries, and inter-node communication grow proportionally. The forensic ledger's append-only property implies monotonic growth without compaction: no garbage collection or archival policy is specified, and no partitioning scheme is defined for distributed deployments.

For deployments at the scale envisioned for Irrig8 across large agricultural regions --- thousands of edge nodes across tens of thousands of acres --- this growth is a practical operational concern. Cross-hierarchy coordination (Corollary~\ref{corol:compositional-safety}) requires explicit Purpose Layer authorization recorded in each participating hierarchy's ledger; as independent hierarchies multiply, coordination overhead at a single anchor creates a bottleneck the protocol does not address.

The 1:1 anchoring model is a structural proof-of-concept. Enterprise and infrastructure deployment requires at minimum 1:N pooled supervision, and potentially N:M anchor pools with load balancing and conflict-resolution protocols, before the protocol can be certified at scale.

\subsubsection{Exploitation Risk from Misconfigured Bounds}
\label{sec:lim-exploitation}

The Bounds Engine enforces depth limits and scope predicates as configured parameters. If $d_{\max}$ is set too large, worst-case escalation latency grows proportionally with hierarchy depth. If scope predicates are miscalibrated --- too permissive --- the bailout triggers so rarely that the protocol becomes a formality; too conservative, and alert fatigue produces operator habituation, causing genuine bailout events to be misclassified as noise.

More insidiously: a sophisticated agent that understands threshold calibration logic may strategically position itself near but not above the confidence threshold $\tau$, exploiting the gap between ``I am uncertain'' and ``I am certain I should escalate.'' This is not a protocol failure; it is a human factors and configuration problem. The defense is procedural: threshold calibration must be subject to human review on a defined schedule, not set once at provisioning and forgotten. Threshold drift detection --- monitoring the distribution of confidence scores over time to detect deliberate positioning --- is an open design question.

% ---------------------------------------------------------------
\subsection{Future Work}
\label{sec:future-work}

Three directions follow directly from the limitations above, presented in order of increasing research complexity.

\subsubsection{Formal Verification of Core Theorems}
\label{sec:fw-formal-verification}

The correctness properties in Section~\ref{sec:correctness} are supported by proof sketches. Full formal verification using a proof assistant (Coq, Lean, or Isabelle/HOL) or a formal specification language (TLA$^+$) would close gaps in the informal arguments, confirm the logical consistency of the three-layer architecture, and produce a machine-checkable artifact suitable for regulatory review.

Particularly valuable: a formal model of the immutable parent pointer $\alpha$ as a hardware-protected function; a formal proof that the forensic ledger's append-only property is preserved under concurrent writes from multiple Fact Layer processes in a distributed deployment; and a compositional verification argument showing that if each layer satisfies its specification, the composed system satisfies the protocol's guarantees. A compositional argument distributes the verification burden across layer implementations and makes the protocol robust to future implementation changes.

\begin{enumerate}
  \item[(i)] \textbf{Fact Layer enforcement.} Formal specification of the pre-commit hook as a state-machine transition guard, proving that no state transition can occur without a corresponding ledger entry and that the ledger entry's issuer field cannot be set to a machine actor.
  \item[(ii)] \textbf{Immutable parent pointer.} Formal proof that $\alpha(n)$ cannot be modified by any machine actor after node provisioning, even under conditions of partial node compromise, subject to the hardware-fault assumption.
  \item[(iii)] \textbf{Compositional safety.} TLA$^+$ specification of the composed three-layer system demonstrating that SAFETY, LIVENESS, and ACCOUNTABILITY hold at the system level as consequences of layer-level invariants.
\end{enumerate}

Formal verification is a prerequisite for deployment in regulated industries where algorithmic accountability is subject to statutory requirements, including EU AI Act compliance for high-risk AI systems.

\subsubsection{Adaptive Timeout and Priority Escalation}
\label{sec:fw-adaptive-timeout}

The fixed timeout model is a conservative baseline. Adaptive timeout mechanisms should adjust the escalation window based on: current depth (nodes near $d_{\max}$ warrant higher urgency than nodes at depth 1); node criticality classification (physical actuator nodes such as irrigation valves warrant faster escalation than informational nodes); historical anchor response patterns (a consistently responsive anchor may be assigned shorter windows, with automatic fallback to a secondary anchor); and environmental conditions (hazardous sensor readings may trigger preemptive escalation before timeout, trading false-positive cost against catastrophic failure cost).

Adaptation requires a feedback loop updating timeout parameters from observed operator response data --- a control theory problem amenable to a proportional-integral-derivative (PID) controller or a lightweight reinforcement learning agent. The critical constraint: adaptation must preserve the deterministic timing guarantee required by the Safety property. The applicable timeout for any given bailout signal must be computable from publicly observable state, preventing the timeout value itself from becoming an attack surface.

The design space includes: a \emph{timeout predictor} that estimates human response time given current depth, node criticality, and time-of-day; a \emph{priority accelerator} that increases $\pi(n, x)$ dynamically when environmental conditions warrant; and a \emph{degradation detector} that identifies anchors whose response patterns have changed significantly and adjusts their $T_{\text{human}}$ allocations accordingly. Each mechanism must be formally specified and its interaction with the Safety and Liveness properties must be proven before deployment.

\subsubsection{Distributed Human Anchor Pools}
\label{sec:fw-anchor-pools}

The 1:N and N:M anchor pool generalizations address the scalability limitation directly. In a 1:N pool, a set $\{h_1, h_2, \dots, h_N\}$ of human anchors share supervision of a single hierarchy, with the Bailout Protocol routing to the first available anchor within $T_{\max}$. In an N:M pool, anchors are partitioned by domain of expertise and jurisdictional scope, requiring the routing algorithm to match exception type to the appropriate anchor's authorized scope.

Critical design constraints for anchor pools:

\begin{enumerate}
  \item[(i)] \textbf{Availability:} escalations must be routed to available anchors within $T_{\max}$ under high concurrent load, requiring load-aware routing rather than static assignment.
  \item[(ii)] \textbf{Accountability:} every escalation must be attributed to a specific human anchor, immutably recorded, requiring that routing decisions are committed to the forensic ledger with the selected anchor's identifier.
  \item[(iii)] \textbf{Cross-jurisdictional authorization:} an anchor in one jurisdiction may lack legal authority to adjudicate actions in another, requiring the accountability chain to record which anchor resolved an escalation and the jurisdictional scope of their authorization.
  \item[(iv)] \textbf{Collusion risk:} pool members may collude to authorize depth budgets that collectively violate safety properties, requiring that anchor pool decisions are logged with enough context for independent audit.
\end{enumerate}

Institutional actors --- corporate compliance teams or designated officers --- introduce a further dimension: institutional rather than individual operational protocols, with chain-of-command rules for anchor delegation. The Treatment of institutional anchors as pool members is a sociotechnical design problem, not solely a protocol design problem, and requires input from legal and organizational scholars.

% ---------------------------------------------------------------
\section{Conclusion}
\label{sec:conclusion}
% ---------------------------------------------------------------

The Bailout Protocol advances a single, architecturally precise claim: that autonomous systems built on the \bxthree{} Framework must propagate every unresolved exception state to a named human principal, bypassing all intermediate machine actors, as a compulsory structural requirement rather than a runtime choice. This mandatory bailout property is not a feature that can be retrofitted to an existing agentic architecture through better error handling or improved escalation heuristics. It is a structural property that follows, as a matter of architectural necessity, from the functional separation between the Purpose, Bounds, and Fact layers.

When the Bounds Layer encounters a condition outside its provisioned operating range that it cannot resolve, no machine actor holds the authority to adjudicate that condition. The decision must return to the Purpose Layer, which carries intent, judgment, and final authority. The Bailout Protocol is the mechanism that enforces this return --- compulsorily, deterministically, and with full forensic attribution at every step.

The contributions of this paper are threefold:

\begin{enumerate}
  \itemsep0em
  \item \textbf{Mandatory architectural bailout protocol.} Recursive Spawning requires every spawned node to maintain an immutable parent pointer whose chain terminates at a human Purpose Layer actor. This is not a best practice or a configuration option --- it is a structural requirement of the spawning mechanism. The Bailout Protocol is the mandatory operationalization of this requirement: it is not invoked as a fallback when something goes wrong, but as the prescribed resolution path for any unresolvable state at any depth.

  \item \textbf{Formal correctness properties.} The drift prevention theorem (Theorem~\ref{thrm:drift-prevention}), chain integrity theorem (Theorem~\ref{thrm:chain-integrity}), and termination theorem (Theorem~\ref{thrm:termination}) establish that the mechanism achieves its stated goals under clearly stated assumptions, providing the analytical foundation for certification, regulatory review, and safety case construction.

  \item \textbf{Integration with the \bxthree{} Framework.} The three-layer architecture is shown to be structurally necessary: the Purpose Layer provides the accountability terminus, the Bounds Engine enforces finite hierarchy depth, and the Fact Layer provides the forensic ledger that makes chain integrity verifiable. Removing any layer renders the mechanism incomplete.
\end{enumerate}

The \bxthree{} Framework is what makes the Bailout Protocol possible. Without a pre-existing three-layer architecture that enforces immutable Purpose encoding at the Fact Layer, an immutable parent pointer is merely a data structure with no enforcement mechanism. The framework provides the enforcement context; the Bailout Protocol provides the structural instantiation of that context applied to hierarchical agent lifecycle management. The upstream accountability guarantee --- that \bxthree{} systems \emph{fail upward into human consciousness, never downward into autonomous chaos} --- and the Bailout Protocol are two aspects of the same architectural commitment.

The limitations identified in Section~\ref{sec:limitations} define the boundary between what the protocol achieves within scope and what must follow in subsequent work. Human response latency, adversarial hardening, ledger scalability, and threshold gaming are constraints on the deployment context, not flaws in the protocol's logical design. They are the research agenda: formal verification to establish correctness proofs rigorously; adaptive timeout to improve operational fit across deployment contexts; and distributed human anchor pools to enable deployment at industrial scale while preserving Safety, Liveness, and Accountability.

The protocol's structural necessity within the \bxthree{} Framework is established by its specification. Its practical necessity in the field is established by the operational reality of autonomous systems that, without it, have no architectural path to human accountability when they encounter conditions beyond their provisioned range. The Bailout Protocol does not eliminate uncertainty from autonomous operation. It eliminates the possibility that uncertainty, propagating through an unconstrained hierarchy, could terminate in a machine actor rather than a human one. That elimination is not a soft property. It is a hard architectural constraint, and it is what the \bxthree{} Framework exists to enforce.
